\chapter{Conclusion}
\label{chap:conclusion}

In conclusion, this project established a fundamental framework for automated quality control of ocular MRI images. By adapting an existing quality control pipeline and developing new anatomy-specific image quality measures (IQMs), a machine learning classifier was trained to effectively distinguish between images of acceptable quality and those of poor quality. The study demonstrated that even with a limited dataset, a model with promising performance could be achieved. A key finding was that IQMs derived from anatomical segmentations were the most powerful predictors of image quality. This underlines the critical importance of accurate segmentation for a robust automated quality control system. This work constitutes a successful proof of concept, laying the foundation for a more comprehensive tool capable of improving the reliability and efficiency of large-scale ocular MRI studies.


% more concis one above
% ----------------------------------------
% This project embarked on the significant challenge of defining quality control criteria for MR-Eye. The developed pipeline extends the FetMRQC framework, a tool originally developed for fetal brain MRI, to establish a robust pipeline for the automated Quality Control (QC) of MR-Eye. The primary objective was to develop a system capable of quantitatively assessing eye MR image quality through a tailored suite of Image Quality Metrics (IQMs) and subsequently to leverage these metrics for training machine learning models for automated quality prediction.

% The project successfully achieved its core aims. Key accomplishments include the meticulous re-targeting of the segmentation data management to accurately process defined ocular anatomical classes (LENS, GLOBE, OPTIC\_NERVE, FAT, MUSCLE), incorporating essential preprocessing steps such as robust label mapping and right-eye spatial filtering. The IQM suite itself was substantially enhanced: existing general-purpose and brain-specific metrics (e.g., CNR, CJV, SNR Dietrich, FBER, EFC) were carefully adapted for ophthalmic application, and novel, eye-specific morphological IQMs – Globe Sphericity and Lens Aspect Ratio – were successfully designed and implemented to capture unique anatomical features. Furthermore, established MRIQC-inspired metrics (QI1, QI2, Tissue-to-Max ratio, RPVE) and comprehensive segmentation-based statistics, including topological features, were integrated, resulting in a rich, quantitative feature set extracted from eye MRI data. Throughout this adaptation, notable technical hurdles pertaining to data type compatibility, argument propagation across modules, handling outputs from external libraries, and ensuring morphological operation versioning were systematically addressed and overcome.

% The adapted MREyeQC system now capably and extracts this extensive set of quantitative features, outputting them into a structured \texttt{IQA.csv} file. This forms the essential input for the machine learning component of the QC pipeline. Preliminary analysis of these IQMs revealed their potential sensitivity to variations in image quality and specific artifacts, as indicated by correlation studies with manual ratings (Section~\ref{fig:feature_importances_results}).

% This project successfully lays a crucial and comprehensive foundation for automated, objective quality assessment in the specialized domain of eye imaging. The developed pipeline and the extensive, tailored IQM set offer a significant step forward from subjective manual QC, paving the way for more reliable and reproducible quantitative studies in ophthalmic research. The systematic approach to adapting a brain-focused framework also provides valuable insights for similar cross-domain QC tool development.

% Future work should prioritize several key areas. Firstly, dataset expansion is paramount, with a focus on acquiring a larger number of eye images and, crucially, enriching the representation of images with diverse quality levels, especially those manually rated as "excluded" or "poor." This will be essential for training more generalizable and robust machine learning models. Secondly, further model exploration and refinement are warranted, including experimenting with advanced feature selection techniques, alternative classification algorithms, and methods specifically designed for imbalanced learning beyond `class\_weight`. Thirdly, the heuristic definitions for complex IQMs such as the artifact indices (QI1, QI2) and RPVE should be further refined and validated, potentially through correlation with specific, visually identified artifacts or through more sophisticated mask generation techniques. Finally, prospective validation of the most promising automated QC model on unseen data from different scanners or sites would be an important step towards demonstrating its broader applicability and potential for clinical or research integration.

% In conclusion, the MREyeQC project has made substantial progress in addressing the need for automated quality control in MR-Eye. While the initial predictive models highlight areas for improvement, the successfully adapted IQM pipeline provides a robust and extensible platform for future advancements, ultimately aiming to enhance the quality and reliability of quantitative eye imaging research and its clinical applications.
